% CP-VI-0044
% target_chapter: VI/41 — Divisor Class Groups
% source_unit_id: IMP-DL-0A9B1F0A75-P001
% source: Downloads/theory-of-algebraic-geometry-9.tex L229-L819
% solution_provenance: SOURCE_EXPLICIT_SOLUTION

\subsection*{Problem statement}

For each of the following schemes, show the claimed isomorphism. In each case, remove a suitable closed subset.

\begin{enumerate}[label=(\alph*)]
	\item
	\[
	\operatorname{Cl}\!\left(\mathbf{P}^1_k \times_{\operatorname{Spec} k} \mathbf{P}^1_k\right)\cong \mathbb{Z}^2.
	\]
	
	\item Let
	\[
	X=\operatorname{Spec} k[x,y,z]/(xy-z^2).
	\]
	Show that
	\[
	\operatorname{Cl}(X)\cong \mathbb{Z}/2\mathbb{Z}.
	\]
	
	\item Let
	\[
	X=\operatorname{Spec} k[x,y,z,w]/(xy-zw).
	\]
	Show that
	\[
	\operatorname{Cl}(X)\cong \mathbb{Z}.
	\]
\end{enumerate}

\subsection{A localization principle for divisor class groups}

We shall repeatedly use the following standard fact.

\begin{proposition}[Removing codimension-one closed subsets]
	Let \(X\) be a normal integral Noetherian scheme, and let
	\[
	U=X\setminus \bigcup_{i=1}^{r}D_i,
	\]
	where \(D_1,\dots,D_r\) are the irreducible codimension-one components of the complement \(X\setminus U\). Then there is an exact sequence
	\[
	\bigoplus_{i=1}^{r}\mathbb{Z}[D_i]
	\longrightarrow
	\operatorname{Cl}(X)
	\longrightarrow
	\operatorname{Cl}(U)
	\longrightarrow 0.
	\]
	Consequently, if \(\operatorname{Cl}(U)=0\), then the divisor classes
	\[
	[D_1],\dots,[D_r]
	\]
	generate \(\operatorname{Cl}(X)\).
\end{proposition}

\begin{proof}
	Every Weil divisor on \(U\) is obtained by restricting a Weil divisor on \(X\). Two divisors on \(X\) restrict to the same divisor on \(U\) precisely when their difference is supported on the codimension-one components of \(X\setminus U\). Passing to divisor classes gives the asserted exact sequence.
\end{proof}

The strategy in all three parts is therefore:

\begin{enumerate}[label=\arabic*.]
	\item Remove suitable divisors so that the remaining open subset is factorial.
	\item Conclude that the removed prime divisors generate the class group.
	\item Compute the relations among their divisor classes using principal divisors.
\end{enumerate}

% ------------------------------------------------------------
\subsection{\texorpdfstring{\(\operatorname{Cl}(\mathbf{P}^1_k\times \mathbf{P}^1_k)\cong \mathbb{Z}^2\)}{Cl(P1 x P1) = Z2}}

\textbf{Solution.}

Let
\[
X=\mathbf{P}^1_k\times_{\operatorname{Spec} k}\mathbf{P}^1_k.
\]

Choose the two divisors
\[
D_1=\{\infty\}\times \mathbf{P}^1_k,
\qquad
D_2=\mathbf{P}^1_k\times \{\infty\}.
\]

Their complement is
\[
U=X\setminus(D_1\cup D_2).
\]

Since
\[
\mathbf{P}^1_k\setminus\{\infty\}\cong \mathbf{A}^1_k,
\]
we obtain
\[
U
\cong
\mathbf{A}^1_k\times \mathbf{A}^1_k
\cong
\mathbf{A}^2_k
=
\operatorname{Spec} k[s,t].
\]

The polynomial ring \(k[s,t]\) is a unique factorization domain. Therefore
\[
\operatorname{Cl}(U)=0.
\]

By the localization principle, the classes of the two removed divisors generate the class group:
\[
\mathbb{Z}[D_1]\oplus \mathbb{Z}[D_2]
\twoheadrightarrow
\operatorname{Cl}(X).
\]

It remains to show that there is no nontrivial relation between \([D_1]\) and \([D_2]\).

Suppose that
\[
aD_1+bD_2=\operatorname{div}(f)
\]
for some rational function
\[
f\in k(X)^\times.
\]

Since the divisor of \(f\) is supported entirely on \(D_1\cup D_2\), the restriction of \(f\) to \(U\) has neither zeros nor poles on \(U\). Thus \(f|_U\) is a unit in
\[
\Gamma(U,\mathcal{O}_U)=k[s,t].
\]

But
\[
k[s,t]^\times=k^\times.
\]

Hence
\[
f|_U=c
\]
for some nonzero constant \(c\in k^\times\). Since \(U\) is dense in \(X\), two rational functions agreeing on \(U\) agree in the function field. Therefore
\[
f=c
\]
as a rational function on \(X\). Consequently,
\[
\operatorname{div}(f)=0.
\]

Thus
\[
aD_1+bD_2=0,
\]
and therefore
\[
a=b=0.
\]

Hence \(D_1\) and \(D_2\) freely generate the divisor class group. Therefore
\[
\boxed{
	\operatorname{Cl}\!\left(\mathbf{P}^1_k\times_{\operatorname{Spec} k}\mathbf{P}^1_k\right)
	\cong
	\mathbb{Z}[D_1]\oplus \mathbb{Z}[D_2]
	\cong
	\mathbb{Z}^2.
}
\]

\begin{remark}
	Geometrically, the two generators correspond to the two different rulings of the surface:
	\[
	\{\text{point}\}\times \mathbf{P}^1_k
	\qquad\text{and}\qquad
	\mathbf{P}^1_k\times\{\text{point}\}.
	\]
	A divisor of the first type cannot be linearly equivalent to a nonzero multiple of a divisor of the second type.
\end{remark}

% ------------------------------------------------------------
\subsection{\texorpdfstring{The quadric cone \(xy=z^2\): \(\operatorname{Cl}(X)\cong \mathbb{Z}/2\mathbb{Z}\)}{The quadric cone xy=z2}}

\textbf{Solution.}

Let
\[
A=k[x,y,z]/(xy-z^2),
\qquad
X=\operatorname{Spec} A.
\]

We shall prove that
\[
\operatorname{Cl}(X)\cong \mathbb{Z}/2\mathbb{Z}.
\]

\subsubsection{Normality of \(X\)}

The scheme \(X\) is the affine quadric cone defined by the equation
\[
xy=z^2.
\]

It is an integral hypersurface. Since hypersurface rings are Cohen--Macaulay, the ring \(A\) satisfies Serre's condition \(S_2\).

The only singular point is the vertex
\[
(x,y,z)=(0,0,0).
\]

Since \(X\) has dimension \(2\), the singular point has codimension \(2\). Thus \(X\) is regular in codimension \(1\), so \(A\) satisfies Serre's condition \(R_1\).

By Serre's normality criterion, \(A\) is normal. Hence the divisor class group \(\operatorname{Cl}(X)\) is defined using Weil divisors.

\subsubsection{Removing a divisor}

Consider the principal open subset
\[
U=D(x)\subseteq X.
\]

On \(U\), the element \(x\) is invertible. From the defining relation
\[
xy=z^2
\]
we obtain
\[
y=\frac{z^2}{x}.
\]

Therefore
\[
A_x
\cong
k[x,x^{-1},z].
\]

The ring \(k[x,x^{-1},z]\) is a Laurent polynomial ring over \(k\), hence it is a unique factorization domain. Therefore
\[
\operatorname{Cl}(U)=0.
\]

Now compute the complement:
\[
X\setminus U=V(x).
\]

Inside \(X\),
\[
A/(x)
\cong
k[y,z]/(z^2).
\]

The unique minimal prime of this quotient is generated by \(z\). Therefore \(V(x)\) has a single irreducible codimension-one component:
\[
D=V(x,z).
\]

By the localization principle, \([D]\) generates \(\operatorname{Cl}(X)\):
\[
\mathbb{Z}[D]\twoheadrightarrow \operatorname{Cl}(X).
\]

\subsubsection{The relation \(2[D]=0\)}

We now compute the principal divisor of \(x\).

Let
\[
\mathfrak p=(x,z)\subseteq A
\]
be the height-one prime corresponding to \(D\). At the generic point of \(D\), the element \(y\) does not belong to \(\mathfrak p\), and is therefore invertible in \(A_{\mathfrak p}\).

From
\[
xy=z^2
\]
we obtain in the local ring \(A_{\mathfrak p}\):
\[
x=\frac{z^2}{y}.
\]

Since \(y\) is a unit and \(z\) generates the maximal ideal of the discrete valuation ring \(A_{\mathfrak p}\), we have
\[
\operatorname{ord}_D(x)=2.
\]

Thus
\[
\operatorname{div}(x)=2D.
\]

Since principal divisors represent zero in the class group,
\[
2[D]=0
\]
in \(\operatorname{Cl}(X)\).

\subsubsection{There are no additional relations}

Suppose that
\[
nD=\operatorname{div}(f)
\]
for some rational function
\[
f\in k(X)^\times.
\]

Since \(\operatorname{div}(f)\) is supported on \(X\setminus U\), the restriction \(f|_U\) is a unit in
\[
A_x\cong k[x,x^{-1},z].
\]

The units of this Laurent polynomial ring are precisely
\[
A_x^\times=k^\times x^{\mathbb{Z}}.
\]

Hence there exist
\[
c\in k^\times,
\qquad
m\in \mathbb{Z}
\]
such that
\[
f|_U=cx^m.
\]

Since \(U\) is dense in \(X\), equality on \(U\) implies equality in the function field:
\[
f=cx^m.
\]

Therefore
\[
\operatorname{div}(f)
=
m\operatorname{div}(x)
=
2mD.
\]

Thus
\[
n=2m.
\]

Hence every relation on \([D]\) is generated by the relation \(2[D]=0\). Therefore
\[
\boxed{
	\operatorname{Cl}\!\left(\operatorname{Spec} k[x,y,z]/(xy-z^2)\right)
	\cong
	\mathbb{Z}[D]/(2[D])
	\cong
	\mathbb{Z}/2\mathbb{Z}.
}
\]

The nonzero class is represented by the ruling line
\[
D=V(x,z).
\]

% ------------------------------------------------------------
\subsection{\texorpdfstring{The quadric cone \(xy=zw\): \(\operatorname{Cl}(X)\cong \mathbb{Z}\)}{The quadric cone xy=zw}}

\textbf{Solution.}

Let
\[
A=k[x,y,z,w]/(xy-zw),
\qquad
X=\operatorname{Spec} A.
\]

We shall prove that
\[
\operatorname{Cl}(X)\cong \mathbb{Z}.
\]

\subsubsection{Normality of \(X\)}

The scheme \(X\) is the affine three-dimensional quadric cone defined by
\[
xy=zw.
\]

Again \(A\) is a hypersurface ring, hence it satisfies Serre's condition \(S_2\).

Its only singular point is the origin
\[
(x,y,z,w)=(0,0,0,0).
\]

Since \(X\) has dimension \(3\), the singular locus has codimension \(3\). In particular, \(X\) is regular in codimension \(1\). Thus \(A\) satisfies \(R_1\).

By Serre's normality criterion, \(A\) is normal.

\subsubsection{Removing a divisor}

Consider
\[
U=D(x)\subseteq X.
\]

Since \(x\) is invertible on \(U\), the equation
\[
xy=zw
\]
gives
\[
y=\frac{zw}{x}.
\]

Therefore
\[
A_x
\cong
k[x,x^{-1},z,w].
\]

This is a Laurent polynomial ring and hence a unique factorization domain. Therefore
\[
\operatorname{Cl}(U)=0.
\]

Now
\[
X\setminus U=V(x).
\]

We compute
\[
A/(x)
\cong
k[y,z,w]/(zw).
\]

The ring \(k[y,z,w]/(zw)\) has two minimal primes:
\[
(z)
\qquad\text{and}\qquad
(w).
\]

Therefore \(V(x)\) has two irreducible codimension-one components:
\[
D=V(x,z),
\qquad
E=V(x,w).
\]

Because \(\operatorname{Cl}(U)=0\), these two divisor classes generate \(\operatorname{Cl}(X)\):
\[
\mathbb{Z}[D]\oplus \mathbb{Z}[E]
\twoheadrightarrow
\operatorname{Cl}(X).
\]

\subsubsection{The relation coming from \(\operatorname{div}(x)\)}

We compute the order of vanishing of \(x\) along \(D\) and \(E\).

At the generic point of
\[
D=V(x,z),
\]
the elements \(y\) and \(w\) are invertible. Since
\[
xy=zw,
\]
we have
\[
x=\frac{w}{y}z.
\]

The factor \(\frac{w}{y}\) is a unit at the generic point of \(D\). Therefore
\[
\operatorname{ord}_D(x)=1.
\]

Similarly, at the generic point of
\[
E=V(x,w),
\]
the elements \(y\) and \(z\) are invertible, and
\[
x=\frac{z}{y}w.
\]

Hence
\[
\operatorname{ord}_E(x)=1.
\]

Therefore
\[
\operatorname{div}(x)=D+E.
\]

Since principal divisors represent zero in the divisor class group,
\[
[D]+[E]=0.
\]

Thus
\[
[E]=-[D],
\]
and \(\operatorname{Cl}(X)\) is generated by the single class \([D]\).

\subsubsection{There is no torsion relation}

Suppose that
\[
aD+bE=\operatorname{div}(f)
\]
for some rational function
\[
f\in k(X)^\times.
\]

Since the divisor of \(f\) is supported on \(X\setminus U\), its restriction to \(U\) is a unit in
\[
A_x\cong k[x,x^{-1},z,w].
\]

The units of this ring are
\[
A_x^\times=k^\times x^{\mathbb{Z}}.
\]

Hence
\[
f|_U=cx^m
\]
for some
\[
c\in k^\times,
\qquad
m\in \mathbb{Z}.
\]

Since \(U\) is dense in \(X\),
\[
f=cx^m
\]
in the function field. Consequently,
\[
\operatorname{div}(f)
=
m\operatorname{div}(x)
=
mD+mE.
\]

Thus every relation among \(D\) and \(E\) is an integral multiple of
\[
D+E.
\]

Therefore
\[
\operatorname{Cl}(X)
\cong
\frac{\mathbb{Z}[D]\oplus \mathbb{Z}[E]}
{\mathbb{Z}([D]+[E])}.
\]

Since the quotient of \(\mathbb{Z}^2\) by the subgroup generated by \((1,1)\) is isomorphic to \(\mathbb{Z}\), we obtain
\[
\boxed{
	\operatorname{Cl}\!\left(\operatorname{Spec} k[x,y,z,w]/(xy-zw)\right)
	\cong
	\mathbb{Z}.
}
\]

A generator may be taken to be
\[
[D]=[V(x,z)],
\]
while
\[
[E]=-[D].
\]

% ============================================================
